% TOP_PROBLEM: {"id": 6500003, "title": "Morse-Sard Questions for Endpoint Maps", "category_id": 6, "category": {"id": 6, "name": "geometry", "display_name": "Geometry", "description": "Euclidean and non-Euclidean geometry, geometric structures.", "slug": "geometry", "order_index": 6, "created_at": "2026-07-31T15:26:25.671Z"}, "status": "open"}
% FIRST_POSTED: null
% ENTRY_KIND: statement_only
% Imported from: batch13.tex; UnsolvedMath ID 6500003
\documentclass[11pt]{article}
\usepackage[margin=1in]{geometry}
\usepackage{amsmath,amssymb,mathtools}
\usepackage{fontspec,unicode-math}
\setmainfont{Latin Modern Roman}
\setsansfont{Latin Modern Sans}
\setmathfont{Latin Modern Math}
\usepackage{xurl,hyperref}
\hypersetup{colorlinks=true,urlcolor=blue,linkcolor=blue}
\newcommand{\sref}[2]{\hyperlink{source:#1:#2}{[S#2]}}
\newcommand{\cataloguescope}[1]{\paragraph{Recorded mathematical question.}#1}
\begin{document}
% UnsolvedMath ID 6500003; source catalogue code AMR-064-0003
\setcounter{section}{6500002}
\section{Morse-Sard Questions for Endpoint Maps}\label{6500003}
{\small\sffamily UnsolvedMath ID 6500003 · Geometry}
\subsection{Definitions and mathematical statement}

\paragraph{Shared notation.}$\mathbb N=\{1,2,\ldots\}$, and $\mathbb Z,\mathbb Q,\mathbb R,\mathbb C$ denote the integers, rationals, reals and complex numbers. Any other conventions are specified locally.


A sub-Riemannian structure consists of a smooth bracket-generating distribution $D\subset TM$ and a smooth inner product on $D$. A horizontal (admissible) curve has $\dot\gamma(t)\in D_{\gamma(t)}$ almost everywhere, and length $L(\gamma)=\int\|\dot\gamma(t)\|_D\,dt$. The distance is the infimum of lengths of horizontal connections. Fix $q_0\in M$ and let $\mathcal H_{q_0}$ be the Hilbert manifold of $H^1$ horizontal paths $\gamma:[0,1]\to M$ starting at $q_0$. Its endpoint map is $\mathcal E(\gamma)=\gamma(1)$. A curve is singular (abnormal) if the differential $D_\gamma\mathcal E$ is not surjective onto $T_{\gamma(1)}M$. It is optimal singular if it is also length minimizing between its endpoints. Positive measure refers to the smooth measure class on $M$, equivalently any smooth positive volume density.
\paragraph{Statement recorded in the supplied source.}
Can the endpoints of all singular curves starting at $q_0$ fill all of $M$? Can the endpoints of optimal singular curves starting at $q_0$ fill a subset of positive measure in $M$? \sref{6500003}{2}
\subsection{Short English statement}
How large can the endpoint sets of abnormal curves and of minimizing abnormal curves be?
\subsection{Sources}
\begin{description}
\item[\hypertarget{source:6500003:1}{S1}] ULAM.AI and the UnsolvedMath Contributors, \emph{UnsolvedMath: A Curated Collection of Open Mathematics Problems}, record 6500003 (AMR-064-0003). CC BY 4.0. \url{https://huggingface.co/datasets/ulamai/UnsolvedMath}.
\item[\hypertarget{source:6500003:2}{S2}] A. A. Agrachev, \emph{Some open problems}, arXiv:1304.2590v2 (2013), Section III, pp. 3--4. \url{https://arxiv.org/pdf/1304.2590}.
\end{description}
\label{end:6500003}
\end{document}
